% GENERATED FILE. Edit canonical lessons, then run npm run generate:books.
% canonical-source-hash: fnv1a64:a9a85a28eba881aa
% canonical-lessons: RU-C06-voda, RU-C06-kofe, RU-C06-chai, RU-C06-khleb

\chapter{Water, Coffee, Tea, and Bread}
\label{ch:russian-water-coffee-tea-bread}
\hlchaptermodality{6}

\begin{chapteropening}
\textbf{By the end of this chapter:} \emph{I can name water, coffee, tea and bread with their genders in Russian, ask for any of them politely, and say which are native and which are loans.}

The chapter closes on \ru{хлеб}, the word that looks native but is a prehistoric Germanic loan: the reader produces all four words with their genders, extends \ru{пожалуйста} into a real request pattern on \ru{кофе} and \ru{чай}, and sorts the set into native (\ru{вода}), loan disguised by age (\ru{хлеб}), and two visibly foreign words that arrived by opposite routes (\ru{кофе} by sea, \ru{чай} overland). It reaches back to Chapter 5's \ru{любить}, closing the chapter's own final atoms which nothing had yet revisited.
\end{chapteropening}

\section[vodá]{\ru{вода} — \textquotedblleft{}water,\textquotedblright{} and how every Russian noun announces itself}

\label{lesson:RU-C06-voda}



\begin{quote}
Five chapters of verbs, and not one noun. Here is the first — and it comes with a rule that will save you for the rest of the book.
\end{quote}

\subsection*{You'll want to know first — \ru{вода}}
\begin{quote}
\textbf{\ru{вода}} — \emph{vodá} — \textbf{water}
\end{quote}

Stress on the last syllable, \emph{vo-DÁ}; the unstressed first \ru{о} drifts toward \emph{a}, the same reduction you have been hearing since \emph{\ru{спасибо}}.

\begin{grammarlens}[title={Grammar Lens: no \textquotedblleft{}the,\textquotedblright{} no \textquotedblleft{}a\textquotedblright{} — and the ending tells you the gender}]
Two facts belong together, because each explains the other.

\textbf{First: Russian has no articles at all.} No word for \emph{the}, no word for \emph{a}. \emph{\ru{Вода}} by itself can mean \textquotedblleft{}water,\textquotedblright{} \textquotedblleft{}the water,\textquotedblright{} or \textquotedblleft{}a water\textquotedblright{} — context carries the difference. Every Romance or Germanic language taught in this program has spent a whole lesson on articles. Russian never will, because it has none to teach.

\textbf{Second: every Russian noun is one of three genders — masculine, feminine, or neuter — and, unlike Hindi's arbitrary assignment, Russian's \emph{ending} usually tells you which:}

\noindent
\begin{tabularx}{\linewidth}{@{}>{\raggedright\arraybackslash}X>{\raggedright\arraybackslash}X>{\raggedright\arraybackslash}X@{}}
\toprule
\textbf{ending} & \textbf{gender} & \textbf{this chapter's example} \\
\midrule
a consonant & masculine & \emph{(next lesson)} \\
\textbf{-\ru{а} / -\ru{я}} & feminine & \textbf{\ru{вод}-\ru{а}} \\
\textbf{-\ru{о} / -\ru{е}} & neuter & \emph{(ahead)} \\
\bottomrule
\end{tabularx}

\emph{\ru{Вода}} ends in \textbf{-\ru{а}}, so it is \textbf{feminine} — regular, no exception, the easiest kind of noun this language has. Hold that confidence for exactly one lesson; the next word breaks the rule on purpose.
\end{grammarlens}

\begin{cousinweb}
\textbf{\ru{вода}} is Indo-European *\emph{wódr}, one of the best-preserved roots in the whole family. English \textbf{water} is the same word, barely disguised; so is the combining form \textbf{hydro-} in \emph{hydrogen} and \emph{hydrate}, which English borrowed separately from the Greek branch of the same root.

Russian kept building on it at home, too: \textbf{vódka} is \emph{\ru{вод}-\ru{а}} plus the affectionate diminutive \textbf{-\ru{ка}} — literally \textbf{\textquotedblleft{}little water.\textquotedblright{}}
\end{cousinweb}

\subsection*{Your turn}
Say these aloud:

\begin{itemize}
\raggedright
  \item \textquotedblleft{}\ru{вода}\textquotedblright{} — \emph{vo-DÁ}, stress on the end
  \item the three endings — \textquotedblleft{}consonant, masculine; -\ru{а} or -\ru{я}, feminine; -\ru{о} or -\ru{е}, neuter\textquotedblright{}
  \item \textquotedblleft{}water, hydro-, vódka\textquotedblright{} — one root, three words you already own
\end{itemize}

\begin{tcolorbox}[breakable,colback=teal!4,colframe=teal!35!black,title={Before you move on}]
Say \textquotedblleft{}water.\textquotedblright{} (\textbf{\ru{Вода}}.) What two words does Russian have for \textquotedblleft{}the\textquotedblright{} and \textquotedblleft{}a\textquotedblright{}? (\textbf{Neither exists.}) What ending makes \emph{\ru{вода}} feminine, and what are the other two endings for? (\textbf{-\ru{а}}; a consonant for masculine, -\ru{о} or -\ru{е} for neuter.) Name two English words from the same root as \emph{\ru{вода}}. (\textbf{Water}, \textbf{hydro-} — and Russian's own \textbf{vódka}, \textquotedblleft{}little water.\textquotedblright{})
\end{tcolorbox}
\section[kófe]{\ru{кофе} — \textquotedblleft{}coffee,\textquotedblright{} the noun that breaks the rule you just learned}

\label{lesson:RU-C06-kofe}



\begin{quote}
Last lesson's rule said -\ru{о} and -\ru{е} endings are neuter. This word ends in -\ru{е} and is not neuter at all — meet Russian's most argued-about noun.
\end{quote}

\subsection*{You'll want to know first — \ru{кофе}}
\begin{quote}
\textbf{\ru{кофе}} — \emph{kófe} — \textbf{coffee}. \textbf{Masculine.}
\end{quote}

Say that twice: \textbf{masculine}, ending in \textbf{-\ru{е}}. By last lesson's rule it should be neuter, like every other -\ru{е} noun ahead in this book. It never declines at all — \emph{\ru{кофе}} is \textbf{indeclinable}, the identical shape in every sentence, so there is no ending left to fix.

\subsection*{The letters in this word}
\emph{(Skim if you read Cyrillic.)} One new letter: \textbf{\ru{ф}}.

\begin{itemize}
\raggedright
  \item \textbf{\ru{ф}} = \textbf{\textquotedblleft{}f\textquotedblright{}} — Greek \textbf{phi} (Φ), a circle with a line through it. You have already met two other Greek-shaped letters, \textbf{\ru{г}} (gamma) and \textbf{\ru{л}} (lambda); \textbf{\ru{ф}} joins them.
\end{itemize}

Everything else in \emph{\ru{кофе}} — \textbf{\ru{к}, \ru{о}, \ru{е}} — you have read since Chapter 1.

\begin{cousinweb}
\textbf{\ru{кофе}} is a loanword, and a recent one beside \emph{\ru{вода}}'s prehistoric root. It travelled \textbf{Arabic \emph{qahwa} $\to$ Ottoman Turkish \emph{kahve} $\to$ Dutch \emph{koffie},} and reached Russian during Peter the Great's westernizing reforms in the early 1700s — the same push that filled Russian with borrowed Dutch and German trade and technical vocabulary.

Its first Russian form was \textbf{\ru{кофий}}, ending like an ordinary masculine noun (compare \emph{\ru{чай}}, next lesson). Across the eighteenth and nineteenth centuries the final syllable wore down to today's \emph{\ru{кофе}} — but the \textbf{gender stuck}, long after the ending that explained it was gone. Dictionaries still mark it masculine; everyday speech increasingly treats it as neuter, exactly the way last lesson's rule would predict, and the argument between the two is still alive in Russian classrooms today.

You already own the verb for how you feel about it:

\begin{quote}
\textbf{\ru{Я} \ru{люблю} \ru{кофе}.} — \emph{ya lyublyú kófe} — \textbf{I love coffee.}
\end{quote}

\emph{(\ru{Кофе} is inanimate and masculine, so it takes exactly the same shape here as it does standing alone — nothing to adjust.)}
\end{cousinweb}

\subsection*{Your turn}
Say these aloud:

\begin{itemize}
\raggedright
  \item \textquotedblleft{}\ru{кофе}\textquotedblright{} — masculine, indeclinable, never changes shape
  \item the new letter — \textquotedblleft{}\ru{ф} is phi, and says f\textquotedblright{}
  \item the route — \textquotedblleft{}qahwa, kahve, koffie, \ru{кофе}\textquotedblright{}
  \item \textquotedblleft{}\ru{Я} \ru{люблю} \ru{кофе}.\textquotedblright{}
\end{itemize}

\begin{tcolorbox}[breakable,colback=teal!4,colframe=teal!35!black,title={Before you move on}]
Say \textquotedblleft{}coffee,\textquotedblright{} with its gender. (\textbf{\ru{Кофе}} — \textbf{masculine}.) Why is that surprising, given last lesson's rule? (\textbf{It ends in -\ru{е}}, which the rule predicts as neuter.) Name the three languages the word crossed before Russian, and the century it arrived. (\textbf{Arabic, Turkish, Dutch}; \textbf{the early 1700s}.) Say \textquotedblleft{}I love coffee.\textquotedblright{} (\textbf{\ru{Я} \ru{люблю} \ru{кофе}}.)
\end{tcolorbox}
\section[chai]{\ru{чай} — \textquotedblleft{}tea,\textquotedblright{} which went overland while \ru{кофе} went by sea}

\label{lesson:RU-C06-chai}



\begin{quote}
Same drink family, a completely different journey — and, finally, the sentence that actually orders one.
\end{quote}

\subsection*{You'll want to know first — \ru{чай}}
\begin{quote}
\textbf{\ru{чай}} — \emph{chai} — \textbf{tea}. \textbf{Masculine.}
\end{quote}

A consonant ending, so last lesson's rule holds without a fight this time: consonant-final nouns are masculine, and \emph{\ru{чай}} is the plain case. Every letter — \textbf{\ru{ч}, \ru{а}, \ru{й}} — you have read since Chapter 2.

\emph{(Careful on paper: romanized, this word looks like the English drink brand name \textbf{chai} and nothing else — pure coincidence of spelling, not a shared history. Its real cousins are below.)}

\begin{grammarlens}[title={Grammar Lens: asking politely, at last}]
You have owned \textbf{\ru{пожалуйста}} since Chapter 1 and had nothing new to attach it to. Now you do: name the thing, add the word.

\begin{quote}
\textbf{\ru{Чай}, \ru{пожалуйста}.} — \emph{chai, pazhálusta} — \textbf{Tea, please.} \textbf{\ru{Кофе}, \ru{пожалуйста}.} — \emph{kófe, pazhálusta} — \textbf{Coffee, please.}
\end{quote}

Both words sit here exactly as they do alone — \emph{\ru{чай}} is already the shape a request needs, and \emph{\ru{кофе}} never changes shape at all. That is not an accident of this chapter: it is the simplest version of a pattern you will meet in its fuller, case-shifting form once this book reaches nouns that \emph{do} change for a request. For now, name the thing, add \emph{\ru{пожалуйста}}, and you are understood.
\end{grammarlens}

\begin{cousinweb}
\textbf{\ru{чай}} is Mandarin Chinese \textbf{chá}, and it reached Russia by the opposite road from \emph{\ru{кофе}}. Persian traders on the Silk Road bought it from northern China as \emph{cha-ye khatai}, \textquotedblleft{}tea of Cathay,\textquotedblright{} and their Turkic-speaking customers misheard the trailing syllable as part of the word itself — giving \textbf{Turkish \emph{çay}, Persian \emph{chāy},} and, at the end of the same overland chain, \textbf{Russian \emph{\ru{чай}}}.

Compare last lesson's \emph{\ru{кофе}}: \textbf{Arabic $\to$ Turkish $\to$ Dutch, by ship}, arriving in Peter the Great's Russia from the west. \emph{\ru{Чай}} came \textbf{overland, from the east}, along a caravan route that also carried it into Hindi (\emph{chai}) and Arabic (\emph{shai}). One plant; two words for it in English (\textbf{tea} took the same Chinese root by sea, through Dutch traders on the coast) — and Russian happens to have picked the \textbf{eastern} branch, while English picked the \textbf{southern coastal} one.
\end{cousinweb}

\subsection*{Your turn}
Say these aloud:

\begin{itemize}
\raggedright
  \item \textquotedblleft{}\ru{чай}\textquotedblright{} — one syllable, masculine
  \item \textquotedblleft{}\ru{Чай}, \ru{пожалуйста}. \ru{Кофе}, \ru{пожалуйста}.\textquotedblright{}
  \item the fork — \textquotedblleft{}chá overland to chai; chá by sea to tea\textquotedblright{}
\end{itemize}

\begin{tcolorbox}[breakable,colback=teal!4,colframe=teal!35!black,title={Before you move on}]
Say \textquotedblleft{}tea,\textquotedblright{} with its gender. (\textbf{\ru{Чай}} — \textbf{masculine}.) Ask for it politely. (\textbf{\ru{Чай}, \ru{пожалуйста}.}) Which route did \emph{\ru{чай}} travel, and which did \emph{\ru{кофе}} travel? (\textbf{\ru{Чай}}: overland, Chinese $\to$ Persian $\to$ Turkic $\to$ Russian. \textbf{\ru{Кофе}}: by sea, Arabic $\to$ Turkish $\to$ Dutch $\to$ Russian.) Which English word took the same Chinese root as \emph{\ru{чай}}, but by the other route? (\textbf{Tea} — by sea, through Dutch.)
\end{tcolorbox}
\section[khleb]{\ru{хлеб} — \textquotedblleft{}bread,\textquotedblright{} the chapter's fourth word and its only immigrant besides \ru{кофе}}

\label{lesson:RU-C06-khleb}



\begin{quote}
One more request word, one more letter, and then the whole table — four things to ask for, and what each one's history actually is.
\end{quote}

\subsection*{You'll want to know first — \ru{хлеб}}
\begin{quote}
\textbf{\ru{хлеб}} — \emph{khleb} — \textbf{bread}. \textbf{Masculine.}
\end{quote}

Consonant-final, so the ending rule holds again. And the pattern from last lesson still works:

\begin{quote}
\textbf{\ru{Хлеб}, \ru{пожалуйста}.} — \emph{khleb, pazhálusta} — \textbf{Bread, please.}
\end{quote}

\subsection*{The letters in this word}
\emph{(Skim if you read Cyrillic.)} One new letter: \textbf{\ru{х}}.

\begin{itemize}
\raggedright
  \item \textbf{\ru{х}} = a raspy \textbf{\textquotedblleft{}kh,\textquotedblright{}} made where you'd clear your throat — not the hard \emph{k} of \emph{\ru{кофе}}'s \textbf{\ru{к}}. It \textbf{looks like Latin \emph{x}} and says nothing like it: a fourth false friend, alongside \textbf{\ru{в}, \ru{р}, \ru{с}, \ru{н}} from Chapter 1.
\end{itemize}

\textbf{\ru{л}, \ru{е}, \ru{б}} you already have. Keep \textbf{\ru{х}} apart from last lesson's \textbf{\ru{ф}}: one is a throat rasp, the other a clean \emph{f}.

\begin{cousinweb}
Here is the twist this chapter has been building to. \textbf{\ru{хлеб}} looks and sounds completely native — nothing about it announces itself as foreign the way \emph{\ru{кофе}} or \emph{\ru{чай}} do. But it is a \textbf{borrowing}, and an ancient one: Proto-Slavic took it from \textbf{Proto-Germanic *hlaibaz}, the same root that gave English \textbf{loaf} and German \textbf{Laib}. It happened so far back in prehistory — long before \emph{\ru{кофе}}'s eighteenth century or \emph{\ru{чай}}'s Silk Road centuries — that the word has had well over a thousand years to stop feeling borrowed at all.

So the chapter's four words sort into a pattern worth naming outright:

\noindent
\begin{tabularx}{\linewidth}{@{}>{\raggedright\arraybackslash}X>{\raggedright\arraybackslash}X>{\raggedright\arraybackslash}X@{}}
\toprule
\textbf{word} & \textbf{route} & \textbf{age} \\
\midrule
\textbf{\ru{вода}} & inherited straight from PIE & prehistoric, native \\
\textbf{\ru{кофе}} & Arabic $\to$ Turkish $\to$ Dutch, \textbf{by sea} & 18th century \\
\textbf{\ru{чай}} & Chinese $\to$ Persian $\to$ Turkic, \textbf{overland} & centuries old \\
\textbf{\ru{хлеб}} & Germanic *hlaibaz & prehistoric, but \textbf{borrowed} \\
\bottomrule
\end{tabularx}

\emph{\ru{Вода}} and \emph{\ru{хлеб}} both feel completely native — but only one of them actually is. Age can disguise a loanword as thoroughly as it can preserve an inheritance.
\end{cousinweb}

\subsection*{Your turn}
Say these aloud:

\begin{itemize}
\raggedright
  \item \textquotedblleft{}\ru{хлеб}\textquotedblright{} — the raspy \ru{х}, masculine
  \item \textquotedblleft{}\ru{Хлеб}, \ru{пожалуйста}.\textquotedblright{}
  \item all four requests — \textquotedblleft{}\ru{Вода}, \ru{кофе}, \ru{чай}, \ru{хлеб}, \ru{пожалуйста}\textquotedblright{} — one request word covering the whole table
  \item \textquotedblleft{}\ru{Я} \ru{люблю} \ru{хлеб}.\textquotedblright{} — the same pattern that closed \ru{кофе}'s lesson
\end{itemize}

\begin{tcolorbox}[breakable,colback=teal!4,colframe=teal!35!black,title={Before you move on}]
Say \textquotedblleft{}bread,\textquotedblright{} with its gender, and ask for it politely. (\textbf{\ru{Хлеб}} — \textbf{masculine}; \textbf{\ru{Хлеб}, \ru{пожалуйста}}.) What does \textbf{\ru{х}} sound like, and which Latin letter does it impersonate? (A raspy \textbf{kh}; looks like \textbf{x}.) Which of this chapter's four words is native, which is borrowed but looks native, and which two are visibly foreign? (\textbf{\ru{Вода}} native; \textbf{\ru{хлеб}} borrowed from Germanic but disguised by age; \textbf{\ru{кофе}} and \textbf{\ru{чай}}, both visibly loans, by opposite routes.) Say \textquotedblleft{}I love bread,\textquotedblright{} recalling the verb Chapter 5 ended on. (\textbf{\ru{Я} \ru{люблю} \ru{хлеб}} — \emph{\ru{люблю}}, with its inserted \textbf{\ru{л}}.)
\end{tcolorbox}
